\documentclass[10pt,a4paper]{amsart}
\usepackage{amsmath,amssymb,amsthm,graphicx}
% For commutative diagrams we use Xy-pic
\usepackage[all]{xy}

% We use 2cell for 2-commutative diagrams.
\xyoption{2cell}
% Theorem environments.
%
\theoremstyle{plain}
\newtheorem{theorem}[subsection]{Theorem}
\newtheorem{proposition}[subsection]{Proposition}
\newtheorem{lemma}[subsection]{Lemma}

\theoremstyle{definition}
\newtheorem{definition}[subsection]{Definition}
\newtheorem{example}[subsection]{Example}
\newtheorem{exercise}[subsection]{Exercise}
\newtheorem{situation}[subsection]{Situation}

\theoremstyle{remark}
\newtheorem{remark}[subsection]{Remark}
\newtheorem{remarks}[subsection]{Remarks}

\numberwithin{equation}{subsection}

% Macros
%
\def\lim{\mathop{\mathrm{lim}}\nolimits}
\def\colim{\mathop{\mathrm{colim}}\nolimits}
\def\Spec{\mathop{\mathrm{Spec}}}
\def\Hom{\mathop{\mathrm{Hom}}\nolimits}
\def\Ext{\mathop{\mathrm{Ext}}\nolimits}
\def\SheafHom{\mathop{\mathcal{H}\!\mathit{om}}\nolimits}
\def\SheafExt{\mathop{\mathcal{E}\!\mathit{xt}}\nolimits}
\def\Sch{\mathit{Sch}}
\def\Mor{\mathop{\mathrm{Mor}}\nolimits}
\def\Ob{\mathop{\mathrm{Ob}}\nolimits}
\def\Sh{\mathop{\mathit{Sh}}\nolimits}
\def\NL{\mathop{N\!L}\nolimits}
\def\CH{\mathop{\mathrm{CH}}\nolimits}
\def\proetale{{pro\text{-}\acute{e}tale}}
\def\etale{{\acute{e}tale}}
\def\QCoh{\mathit{QCoh}}
\def\Ker{\mathop{\mathrm{Ker}}}
\def\Im{\mathop{\mathrm{Im}}}
\def\Coker{\mathop{\mathrm{Coker}}}
\def\Coim{\mathop{\mathrm{Coim}}}

% Boxtimes
%
\DeclareMathSymbol{\boxtimes}{\mathbin}{AMSa}{"02}

%
% Macros for moduli stacks/spaces
%
\def\QCohstack{\mathcal{QC}\!\mathit{oh}}
\def\Cohstack{\mathcal{C}\!\mathit{oh}}
\def\Spacesstack{\mathcal{S}\!\mathit{paces}}
\def\Quotfunctor{\mathrm{Quot}}
\def\Hilbfunctor{\mathrm{Hilb}}
\def\Curvesstack{\mathcal{C}\!\mathit{urves}}
\def\Polarizedstack{\mathcal{P}\!\mathit{olarized}}
\def\Complexesstack{\mathcal{C}\!\mathit{omplexes}}
% \Pic is the operator that assigns to X its picard group, usage \Pic(X)
% \Picardstack_{X/B} denotes the Picard stack of X over B
% \Picardfunctor_{X/B} denotes the Picard functor of X over B
\def\Pic{\mathop{\mathrm{Pic}}\nolimits}
\def\Picardstack{\mathcal{P}\!\mathit{ic}}
\def\Picardfunctor{\mathrm{Pic}}
\def\Deformationcategory{\mathcal{D}\!\mathit{ef}}
\newenvironment{reference}{}{}
\newenvironment{slogan}{}{}
\newenvironment{history}{}{}
\usepackage{cleveref}
\crefname{originallemma}{original lemma}{original lemmas}
\Crefname{originallemma}{Original lemma}{Original lemmas}
\makeatletter
\newlabel{lemma-non-irreducible-minimal-type-genus-at-least-one}{{0C7B}{source}{}{sourcepointer-1}{}}
\newlabel{lemma-non-irreducible-minimal-type-genus-at-least-one@cref}{{[originallemma][0][]0C7B}{[1][0][]source}{}{}{}}
\makeatother
\title{A minimal numerical type with more than one vertex has genus at least the topological genus}
\author{The Stacks Project Authors}
\date{stacks-project:models.tex:local-snapshot-20261012; GNU FDL 1.2 or later}
\begin{document}
\maketitle
\noindent\textit{Editorial note.} The counted unit is the Stacks Project Lemma lemma-genus-nonnegative at source lines 911-915 together with its proof at source lines 917-1060, a straight adjacent pairing in the chapter file and the longest proof in this chapter. The chapter is the upstream file models.tex; the pool directory carries the longer name stacks-models and is not used as the chapter name anywhere in this card.\par\medskip
\noindent\textit{Editorial note.} Two context spans carry the definitions the counted statement uses: the lemma that introduces the topological genus at lines 861-866 and the definition of a minimal numerical type at lines 887-893. Both are complete environments from their opening to their closing marker, both are context rather than counted units, and neither overlaps the counted interval.\par\medskip
\noindent\textit{Editorial note.} The retained proof names one lemma of the same chapter at source lines 988 and 1041, the description of minimal types that are not irreducible; it is declared as a non-load-bearing navigation pointer with its stable Stacks tag. The lemma defining the topological genus is retained as context and therefore must NOT also be declared as a pointer: doing so duplicated the label and the engine refused the file with a multiply-defined-label warning.\par\medskip
\noindent\textit{Editorial note.} The author opens the proof with a sentence telling the reader who only cares about types coming from proper regular models that the proof may be skipped, because the geometric case is reproved later in the chapter. That sentence is part of the author's own argument and is retained verbatim; it is not an omission.\par\medskip
\noindent\textit{Editorial note.} The proof is the author's complete argument and contains no omission wording anywhere in its 141 lines.\par\medskip
\noindent\textit{Editorial note.} No citation occurs in the retained spans, so no bibliography entry is supplied.\par\medskip
\noindent\textit{Allowed standard prerequisites (editorial).}\begin{itemize}
\item Numerical types, the topological genus and minimality. A numerical type is a collection of integers n, m of subscript i, a of subscript i j, w of subscript i and g of subscript i satisfying the symmetry and recurrence conditions of the chapter; its topological genus is the expression one minus the number of vertices plus the number of pairs of distinct vertices whose intersection number is positive, and it is nonnegative; the type is minimal when no index has genus zero and self-intersection equal to the negative of its weight, that is, when there is no minus-one index. The author's own definitions are the allowed inputs.. The counted lemma is the lower bound for the genus of a minimal type that the section uses when it classifies types of low genus and when it compares minimal models with their numerical types.
\end{itemize}\medskip
\ifdefined\hypertarget\hypertarget{sourcepointer-1}{}\fi
\noindent\textit{Source reference.} Original lemma 0C7B: focused-seeds/AGP/sources/stacks-snapshot-20261011/tags.txt, source lines 15326--15326. Navigation pointer invoked by the retained proof at source lines 988 and 1041, where the author recalls the description of a minimal type that is not irreducible; the cited statement is an allowed input of the same section and its proof is not reproduced inside the span. The number is the stable Stacks tag of the target, looked up in the shared tag index.\par\medskip
\begin{lemma}
\label{lemma-top-genus}
Let $n, m_i, a_{ij}, w_i, g_i$ be a numerical type.
Let $e$ be the number of pairs $(i, j)$ with $i < j$ and $a_{ij} > 0$.
Then the expression $g_{top} = 1 - n + e$ is $\geq 0$.
\end{lemma}

\begin{definition}
\label{definition-type-minimal}
We say the numerical type $n, m_i, a_{ij}, w_i, g_i$ of genus $g$
is {\it minimal} if there does not exist an $i$
with $g_i = 0$ and $a_{ii} = -w_i$, in other words, if there
does not exist a $(-1)$-index.
\end{definition}

\begin{lemma}
\label{lemma-genus-nonnegative}
If $n, m_i, a_{ij}, w_i, g_i$ is a minimal numerical type
with $n > 1$, then $g \geq g_{top}$.
\end{lemma}

\begin{proof}
The reader who is only interested in the case of numerical types
associated to proper regular models can skip this proof as we will
reprove this in the geometric situation later.
We can write
$$
g_{top} = 1 - n + \frac{1}{2}\sum\nolimits_{a_{ij} > 0} 1 =
1 + \sum\nolimits_i (-1 +
\frac{1}{2}\sum\nolimits_{j \not = i,\ a_{ij} > 0} 1) 
$$
On the other hand, we have
\begin{align*}
g & =
1 + \sum m_i(w_i(g_i - 1) - \frac{1}{2} a_{ii}) \\
& =
1 + \sum m_iw_ig_i - \sum m_iw_i +
\frac{1}{2} \sum\nolimits_{i \not = j} a_{ij}m_j \\
& =
1 + \sum\nolimits_i
m_iw_i(-1 + g_i + \frac{1}{2} \sum\nolimits_{j \not = i} \frac{a_{ij}}{w_i})
\end{align*}
The first equality is the definition, the second equality uses that
$\sum a_{ij}m_j = 0$, and the last equality uses that
uses $a_{ij} = a_{ji}$ and switching order of
summation. Comparing with the formula for $g_{top}$ we conclude
that the lemma holds if
$$
\Psi_i =
m_iw_i(-1 + g_i + \frac{1}{2} \sum\nolimits_{j \not = i} \frac{a_{ij}}{w_i})
- (-1 + \frac{1}{2}\sum\nolimits_{j \not = i,\ a_{ij} > 0} 1)
$$
is $\geq 0$ for each $i$. However, this may not be the case.
Let us analyze for which indices we can have $\Psi_i < 0$.
First, observe that
$$
(-1 + g_i + \frac{1}{2}\sum\nolimits_{j \not = i} \frac{a_{ij}}{w_i}) \geq
(-1 + \frac{1}{2}\sum\nolimits_{j \not = i,\ a_{ij} > 0} 1)
$$
because $a_{ij}/w_i$ is a nonnegative integer. Since $m_iw_i$ is
a positive integer we conclude that $\Psi_i \geq 0$ as soon as
either $m_iw_i = 1$ or the left hand side of the inequality is $\geq 0$
which happens if $g_i > 0$, or $a_{ij} > 0$ for at least two indices $j$, or
if there is a $j$ with $a_{ij} > w_i$. Thus
$$
P = \{i : \Psi_i < 0\}
$$
is the set of indices $i$ such that $m_iw_i > 1$, $g_i = 0$,
$a_{ij} > 0$ for a unique $j$, and $a_{ij} = w_i$ for this $j$.
Moreover
$$
i \in P \Rightarrow \Psi_i = \frac{1}{2}(-m_iw_i + 1)
$$
The strategy of proof is to show that given $i \in P$ we can borrow a bit
from $\Psi_j$ where $j$ is the neighbour of $i$, i.e., $a_{ij} > 0$.
However, this won't quite work because $j$ may be an index with $\Psi_j = 0$.

\medskip\noindent
Consider the set
$$
Z = \{j : g_j = 0\text{ and }
j\text{ has exactly two neighbours }i, k\text{ with }
a_{ij} = w_j = a_{jk}\}
$$
For $j \in Z$ we have $\Psi_j = 0$. We will consider sequences
$M = (i, j_1, \ldots, j_s)$ where $s \geq 0$,
$i \in P$, $j_1, \ldots, j_s \in Z$, and
$a_{ij_1} > 0, a_{j_1j_2} > 0, \ldots, a_{j_{s - 1}j_s} > 0$.
If our numerical type consists of two indices which are in $P$
or more generally if our numerical type consists of
two indices which are in $P$ and all other indices in $Z$, then
$g_{top} = 0$ and we win by
Lemma \ref{lemma-non-irreducible-minimal-type-genus-at-least-one}.
We may and do discard these cases.

\medskip\noindent
Let $M = (i, j_1, \ldots, j_s)$ be a maximal sequence and let
$k$ be the second neighbour of $j_s$. (If $s = 0$, then $k$
is the unique neighbour of $i$.) By maximality $k \not \in Z$
and by what we just said $k \not \in P$. Observe that
$w_i = a_{ij_1} = w_{j_1} = a_{j_1j_2} = \ldots = w_{j_s} =
a_{j_sk}$. Looking at the definition
of a numerical type we see that
\begin{align*}
m_ia_{ii} + m_{j_1}w_i & = 0,\\
m_iw_i + m_{j_1}a_{j_1j_1} + m_{j_2}w_i & = 0,\\
\ldots & \ldots \\
m_{j_{s - 1}}w_i + m_{j_s}a_{j_sj_s} + m_kw_i & = 0
\end{align*}
The first equality implies $m_{j_1} \geq 2m_i$ because the
numerical type is minimal. Then the second equality implies
$m_{j_2} \geq 3m_i$, and so on. In any case, we conclude that
$m_k \geq 2m_i$ (including when $s = 0$).

\medskip\noindent
Let $k$ be an index such that we have a $t > 0$ and pairwise distinct
maximal sequences $M_1, \ldots, M_t$ as above, with
$M_b = (i_b, j_{b, 1}, \ldots, j_{b, s_b})$
such that $k$ is a neighbour of $j_{b, s_b}$ for $b = 1, \ldots, t$.
We will show that $\Phi_j + \sum_{b = 1, \ldots, t} \Phi_{i_b} \geq 0$.
This will finish the proof of the lemma by what we said above.
Let $M$ be the union of the indices occurring in $M_b$, $b = 1, \ldots, t$.
We write
$$
\Psi_k =
-\sum\nolimits_{b = 1, \ldots, t} \Psi_{i_b} + \Psi_k'
$$
where
\begin{align*}
\Psi_k' & =
m_kw_k\left(-1 + g_k +
\frac{1}{2} \sum\nolimits_{b = 1, \ldots t}
(\frac{a_{kj_{b, s_b}}}{w_k} - \frac{m_{i_b}w_{i_b}}{m_kw_k}) +
\frac{1}{2} \sum\nolimits_{l \not = k,\ l \not \in M}
\frac{a_{kl}}{w_k}
\right) \\
&
-\left(
-1 + \frac{1}{2}\sum\nolimits_{l \not = k,\ l \not \in M,\ a_{kl} > 0} 1
\right)
\end{align*}
Assume $\Psi_k' < 0$ to get a contradiction.
If the set $\{l : l \not = k,\ l \not \in M,\ a_{kl} > 0\}$ is empty,
then $\{1, \ldots, n\} = M \cup \{k\}$ and $g_{top} = 0$
because $e = n - 1$ in this case and the result holds by
Lemma \ref{lemma-non-irreducible-minimal-type-genus-at-least-one}.
Thus we may assume there is at least one such $l$ which contributes
$(1/2)a_{kl}/w_k \geq 1/2$ to the sum inside the first brackets.
For each $b = 1, \ldots, t$ we have
$$
\frac{a_{kj_{b, s_b}}}{w_k} - \frac{m_{i_b}w_{i_b}}{m_kw_k} =
\frac{w_{i_b}}{w_k}(1 - \frac{m_{i_b}}{m_k})
$$
This expression is $\geq \frac{1}{2}$ because $m_k \geq 2m_{i_b}$
by the previous paragraph and is $\geq 1$ if $w_k < w_{i_b}$.
It follows that $\Psi_k' < 0$ implies $g_k = 0$.
If $t \geq 2$ or $t = 1$ and $w_k < w_{i_1}$, then $\Psi_k' \geq 0$
(here we use the existence of an $l$ as shown above) which
is a contradiction too.
Thus $t = 1$ and $w_k = w_{i_1}$. If there at least two nonzero terms
in the sum over $l$ or if there is one such $k$ and $a_{kl} > w_k$, then
$\Psi_k' \geq 0$ as well. The final possibility is that $t = 1$ and
there is one $l$ with $a_{kl} = w_k$. This is disallowed as this would
mean $k \in Z$ contradicting the maximality of $M_1$.
\end{proof}

\end{document}
